\chapter[Quasi-coherent sheaves on preschemes]{Application to quasi-coherent sheaves on preschemes} \label{II}

\numberwithin{equation}{section}

\begin{supproposition} \label{II.1}
Let $X$ be a prescheme, $Z$ a locally closed subset of the form $Z=U-V$, where $U$ and $V$ are two open subsets of $X$ such that $V\subset U$ and that the canonical immersions $U\to X$, $V\to X$ are quasi-compact. Then for every quasi-coherent Module $F$ on $X$, the sheaves ${\SheafH}_Z^i(F)$ are quasi-coherent.
\end{supproposition}

By (\Ref{I}~\Ref{eq:I.24}), there exists an exact sequence of relative cohomology
\[
\to {\SheafH}_U^{i-1}(F) \to {\SheafH}_V^{i}(F) \to {\SheafH}_Z^{i}(F) \to {\SheafH}_U^{i}(F) \to {\SheafH}_V^{i+1}(F) \to.
\]

By (\EGA III~1.4.17), for the ${\SheafH}_Z^i(F)$ to be quasi-coherent it therefore suffices that the ${\SheafH}_U^i(F)$ and the ${\SheafH}_V^i(F)$ be so. One may therefore assume $Z$ open and the canonical immersion $j:Z\to X$ quasi-compact.

Since $Z$ is open one has (\Ref{I}~\Ref{eq:I.22}) a canonical isomorphism:
\[
{\SheafH}_Z^{i}(F) \simeq \R^ij_{*}(F_{|Z})
\]
but $j$ is separated (\EGA I~5.5.1) and quasi-compact, hence (\EGA III~1.4.10) the $\R^ij_{*}(F|Z)={\SheafH}_Z^{i}(F)$ are quasi-coherent, which completes the proof.

\begin{supcorollary} \label{II.2}
Let $Z$ be a closed subset of $X$ such that the canonical immersion \hbox{$X-Z\to X$} is quasi-compact, then the Modules ${\SheafH}_Z^i(F)$ are quasi-coherent.
\end{supcorollary}

\begin{supcorollary} \label{II.3}
If $X$ is locally noetherian, then for every locally closed subset $Z$ of $X$, and every Module $F$ quasi-coherent on $X$, the ${\SheafH}_Z^i(F)$ are quasi-coherent.
\end{supcorollary}

This follows immediately from corollary \Ref{II.2} and from (\EGA I~6.6.4).

\begin{supcorollary} \label{II.4}
Suppose that $X$ is the spectrum of a ring $A$ and let $U$ be a quasi-compact open subset of $X$, $Y=X-U$, $F$ a quasi-coherent Module on $X$; there exists an isomorphism of cohomological functors in $F$:
\begin{equation} \label{eq:II.4.1} {\SheafH}_Y^i(F) = \widetilde{(\H_Y^i(X, F))}.
\end{equation}

One has moreover an exact sequence functorial in $F$:
\begin{equation} \label{eq:II.4.2}
0 \to \H_Y^0(X, F) \to \H^0(X, F) \to \H^0(U, F) \to \H_Y^1(X, F) \to 0
\end{equation}
and isomorphisms functorial in $F$:
\begin{equation} \label{eq:II.4.3}
\H_Y^i(X, F) \simeq \H^{i-1}(U, F), \quad i\geq 2.
\end{equation}
\end{supcorollary}

By corollary~\Ref{II.2}, the ${\SheafH}_Y^i(F)$ are quasi-coherent, since $X$ is affine one therefore has $\H^p(X, {\SheafH}_Y^i(F))=0$ if $p>0$. The spectral sequence (\Ref{I}~\Ref{eq:I.23}) degenerates, hence
\[
\H_Y^i(X, F) = \Gamma({\SheafH}_Y^i(F)).
\]
The equality (\Ref{eq:II.4.1}) then follows from (\EGA I~1.1.3.7), (\Ref{eq:II.4.2}) and (\Ref{eq:II.4.3}) from the exact cohomology sequence (\Ref{I}~\Ref{eq:I.27}) and from the fact that $\H^i(X, F)=0$ if $i>0$, since $X$ is affine.

With the hypotheses of \Ref{II.4}\nde{for the sake of consistency and clarity, one has numbered in parentheses only the equations.}, $U$ is a finite union of affine open sets $X_f$, one can therefore find an ideal $I$ generated by a finite number of elements $f_{\alpha}$ and defining~$Y$, say $\bbf=(f_{\alpha})$. With the notation of (\EGA III~1)\nde{recall that $H^{\boule}(\bbf, M)$ is the Koszul cohomology $H^{\boule}(\Hom(K_{\boule}(\bbf), M))$ of $\bbf$ (\EGA III~1.1.2) with values in $M$ and that $H^{\boule}((\bbf), M)$ is the limit (\loccit, 1.1.6.5) $\varinjlim_n H^{\boule}(\bbf^n, M)$, the transition morphisms being induced by the natural morphisms $K_{\boule}(\bbf^{n+1})\to K_{\boule}(\bbf^n)$ (\loccit, 1.1.6).} one has:

\begin{supproposition} \label{II.5}
Suppose that $X$ is the spectrum of a ring $A$, let $\bbf=(f_{\alpha})$ be a finite family of elements of $A$, $Y$ the closed subset of $X$ that they define, $M$ an $A$-module, $F$ the sheaf associated with $M$. One then has isomorphisms of $\partial$-functors in $M$:
\begin{equation} \label{eq:II.5.1} \H^i((\bbf), M) \simeq \H_Y^i(X, F).
\end{equation}
\end{supproposition}

\noindent
(We shall also write $\H_J^i(M)=\H_Y^i(X, F)$, if $Y$ is the closed subset of $X=\Spec A$ defined by an ideal $J$ of $A$).

For $i=0$ and $i=1$, one uses the exact sequences (\Ref{eq:II.4.2}) and (\EGA III~1.4.3.2); if $i\geq2$, one uses (\Ref{eq:II.4.3}) and (\EGA III~1.4.3.1). This gives us isomorphisms functorial in $M$. One verifies that, up to a sign depending only on $i$, they are compatible with the boundary operator, whence the existence of the isomorphism of $\partial$-functors (\Ref{eq:II.5.1}).

Let now $X$ be a prescheme, $Y$ a closed subset of $X$ and $f:Y\to X$ the inclusion, $I$ a quasi-coherent ideal defining $Y$ in $X$. Let $F$ be a sheaf on $X$.

One has seen that there exist isomorphisms of $\partial$-functors in $F$
\begin{align}
\label{eq:II.5.*} \tag{$*$} \Ext_{\OX}^i(X;f_{*}f^{-1}(\OX), F) &\to \H_Y^i(X, F) \\
\label{eq:II.5.**} \tag{$**$} \SheafExt_{\OX}^i(f_{*}f^{-1}(\OX), F) &\to \SheafH_Y^i(F).
\end{align}

Let $n, m$ be integers such that $m\geq n\geq 0$, one denotes by $i_{n, m}$ the canonical map: $\Oo_{Y_m}=\OX/I^{m+1} \to\OX/I^{n+1}=\Oo_{Y_n}$, and by $j_n$ the map: $f_{*}f^{-1}(\OX)\to\Oo_{Y_n}$. The $(\Oo_{Y_n}, i_{n, m})$ form a projective system and the $j_n$ are compatible with the $i_{n, m}$.

Applying the functor $\Ext_{\OX}^i(X;\cdot, F)$, one deduces from this a morphism
\[
\varphi': \varinjlim_{n} \Ext_{\OX}^i(X;\Oo_{Y_n}, F) \to \Ext_{\OX}^i(X;f_{*}f^{-1}(\OX), F);
\]
one shows easily that it is a morphism of cohomological functors in $F$. The morphism
\[
\varphi: \varinjlim_{n} \Ext_{\OX}^i(X;\Oo_{Y_n}, F) \to \H_Y^i(X, F),
\]
composite of $\varphi'$ and of (\Ref{eq:II.5.*}), is therefore itself also a morphism of cohomological functors in $F$.

One defines similarly
\[
\underline{\varphi}: \varinjlim_{n} \SheafExt_{\OX}^i(\Oo_{Y_n}, F) \to {\SheafH}_Y^i(F).
\]

One has in view the following theorem:

\begin{suptheorem} \label{II.6} \textup{a)}
Let $X$ be a locally noetherian prescheme, $Y$ a closed subset of $X$ defined by a coherent Ideal $I$, $F$ a quasi-coherent Module. Then $\boldsymbol{\varphi}$ is an isomorphism.

\textup{b)} If $X$ is noetherian $\varphi$ is an isomorphism.
\end{suptheorem}

Theorem \Ref{II.6} will follow from \Ref{II.6}.a) and from the

\begin{suplemma} \label{II.7} If the underlying topological space of $X$ is noetherian and if $\boldsymbol{\varphi}$ is an isomorphism, the same is true of $\varphi$.
\end{suplemma}

One will first prove lemma \Ref{II.7}. One knows that there exists a spectral sequence
\begin{equation} \label{eq:II.7.1}
\H^p(X, {\SheafH}_Y^q(F)) \To \H_Y^{\ast}(X, F).
\end{equation}
On the other hand one has an inductive system of spectral sequences
\begin{equation*} \label{eq:II.7.2n} \tag{7.2$_n$}\stepcounter{equation}
\H^p(X, {\SheafExt}_{\OX}^q(\Oo_{Y_n}, F)) \To \Ext_{\OX}^{\ast}(X;\Oo_{Y_n}, F).
\end{equation*}
It follows from the definition of $\varphi$ and of $\underline{\varphi}$ that these morphisms are associated with a homomorphism $\Phi$ of spectral sequences from the inductive limit of (\Ref{eq:II.7.2n}) into (\Ref{eq:II.7.1}). If the underlying space of $X$ is noetherian, by (God. 4.12.1)\sfootnote{\Cf the first bibliographic reference, at the end of \Exp \Ref{I}.}
\begin{equation*}
\varinjlim_{n} \H^p(X, {\SheafExt}_{\OX}^q(\Oo_{Y_n}, F))\isomto \H^p(X, \varinjlim_{n} {\SheafExt}_{\OX}^q(\Oo_{Y_n}, F)),
\end{equation*}
then $\Phi_2$ can be written as a morphism:
\begin{equation*} \H^p(X, \varinjlim_{n} {\SheafExt}_{\OX}^q(\Oo_{Y_n}, F)) \to \H^p(X, {\SheafH}_Y^q(F))
\end{equation*}
which is none other than the one deduced from $\underline{\varphi}$.

If $\underline{\varphi}$ is an isomorphism the same is therefore true of $\Phi_2$, hence of $\varphi$ by ($\EGA 0_{\textup{III}}$ 11.1.5), lemma~\Ref{II.7} is therefore proved.

One will now prove \Ref{II.6}.a); this is a local question on $X$. By corollary~\Ref{II.4} and (\EGA I~1.3.9 and 1.3.12) one may assume that $X$ is the spectrum of a ring $A$. It therefore suffices to prove that under the hypotheses of theorem~\Ref{II.6}.a), the canonical homomorphism:
\begin{equation} \label{eq:II.7.3} \varinjlim_{n} \Ext_A^i(A/I^n, M) \to \H_Y^i(X, M)
\end{equation}
is an isomorphism.

Let $f_{\alpha}$ be a finite number of elements of $A$ generating $I$, $\bbf=(f_{\alpha})$; then the sequence of ideals $(\bbf^n)$ is decreasing and cofinal with the sequence of the $I^n$, so that (\Ref{eq:II.7.3}) is equivalent to a morphism of $\partial$-functors in $M$:
\begin{equation} \label{eq:II.7.4} \varinjlim_{n} \Ext_A^i(A/(\bbf^n), M) \to \H_Y^i(X, M).
\end{equation}

On the other hand one has canonical isomorphisms:
\begin{equation} \label{eq:II.7.5} \varinjlim_{n} \Hom_A(A/(\bbf^n), M) \simeq \varinjlim_{n} (m\in M\ |\ (\bbf^n)m=0) \simeq \H^0((\bbf), M).
\end{equation}
Since $\varinjlim_{n} \Ext_A^i(A/(\bbf^n), M)$ is a universal $\partial$-functor in $M$, there exists a unique morphism of $\partial$-functors in $M$:
\begin{equation} \label{eq:II.7.6} \varinjlim_{n} \Ext_A^i(A/(\bbf^n), M) \to \H^i((\bbf), M),
\end{equation}
which coincides in degree zero with (\Ref{eq:II.7.5}).

Since the composite of (\Ref{eq:II.7.3}) and of (\Ref{eq:II.5.1}) is a morphism of $\partial$-functors in $M$ which coincides with (\Ref{eq:II.7.6}) in degree $0$, it coincides with (\Ref{eq:II.7.6}) in every degree. Theorem \Ref{II.6}.a) is therefore an immediate consequence of the

\begin{suplemma} \label{II.8}
Let $A$ be a noetherian ring, $I$ an ideal generated by a finite system $\bbf=(f_{\alpha})$ of elements, $M$ an $A$-module. Then the homomorphisms (\Ref{eq:II.7.6}) are isomorphisms.
\end{suplemma}

\begin{suplemma} \label{II.9}
Let $A$ be a ring, $\bbf=(f_{\alpha})$ a finite system of elements of $A$, $I$ the ideal generated by $\bbf$, $i$ an integer $>0$. The following conditions are equivalent:
\begin{enumeratea}
\item
The homomorphism (\Ref{eq:II.7.6}) is an isomorphism for every $M$.
\item
$\H^i((\bbf), M)=0$ for $M$ injective.
\item
The projective system $(\H_i(\bbf^n, A))=\H_{i, n}$ is essentially zero, that is to say: for every $n$, there exists $n'>n$ such that $\H_{i, n'}\to\H_{i, n}$ is zero.
\end{enumeratea}
\end{suplemma}

a) implies b) trivially.

b) implies a); indeed b) implies that $M\mto\H^i((\bbf), M)$ is a universal cohomological functor, (\Ref{eq:II.7.6}) is then a morphism of universal cohomological functors. It is an isomorphism in degree zero, hence in every degree.

c) implies b); indeed if $M$ is injective, one has for every $n$
\begin{equation*} \H^i(\bbf^n, M)=\Hom(\H_i(\bbf^n, A), M) = \Hom(\H_{i, n}, M),
\end{equation*}
c) therefore implies that for every $i$ the inductive system $(\H^i(\bbf^n, M))_{n\in\ZZ}$ is essentially zero, whence b).

b) implies c). Indeed, let $n>0$, and let $j$ be a monomorphism of $\H_{i, n}$ into an injective module $M$. Let $n'\geq n$ and let $j_{n'}\in\Hom(\H_{i, n'}, M)$ be the composite of $j$ and of the transition homomorphism $t_{n', n}:\H_{i, n'}\to\H_{i, n}$. The $j_{n'}$ define an element of $\H^i((\bbf), M)$ which is zero by hypothesis. There therefore exists $n_0$ such that $j_{n'}=0$ if $n'>n_0$. But since $j$ is a monomorphism, $j_{n'}=0$ implies $t_{n', n}=0$, whence the proposition.

\begin{supcorollary} \label{II.10} Suppose that the underlying space of $X=\Spec(A)$ is noetherian. For the preceding conditions to be satisfied for every finite family of elements of $A$ and every $i>0$ (or again: for $i=1$), it is necessary and sufficient that for every injective $A$-module $M$, the sheaf $F$ associated with $M$ be flasque.
\end{supcorollary}

This is necessary: indeed, let $\bbf=(f_{\alpha})$ be a finite system of elements of $A$, $Y$ the closed set defined by $\bbf$ and $U=X-Y$; one then has the exact sequence
\begin{equation*} \H^0(X, F) \to \H^0(U, F) \to \H^1((\bbf), M) \to 0,
\end{equation*}
and by \Ref{II.9}.b, $\H^0(X, F)\to\H^0(U, F)$ is surjective.

This is sufficient by virtue of (\Ref{eq:II.5.1}) and of the fact that for every closed subset $Y$ of $X$ and every flasque sheaf $F$ on $X$, $\H_Y^i(X, F)=0$ for $i>0$.

\begin{suplemma} \label{II.11} Under the hypotheses of lemma~\Ref{II.9}, for every noetherian $A$-module $N$ and for every $i>0$, the projective system $(\H_{i, n}(N))_{n\in\ZZ}$, where $\H_{i, n}(N)=\H_i(\bbf^n, N)$, is essentially zero.
\end{suplemma}

Proof by induction on the number $m$ of elements of $\bbf$.

If $m=1$, $\bbf$ is reduced to a single element, say $f$, $\H_{i, n}(N)$ is zero if $i>1$ and $\H_{1, n}(N)$ is canonically isomorphic to the annihilator $N(n)$ of $f^n$ in $N$, the transition homomorphism $N(n')\to N(n)$, $n'\geq n$, being multiplication by $f^{n'-n}$. The $N(n)$ form an increasing sequence of submodules of $N$, and since $N$ is noetherian there exists $n_0$ such that $N(n)=N(n_0)$ if $n\geq n_0$. Hence all the $N(n)$ are annihilated by $f^{n_0}$ and the transition homomorphisms $N(n')\to N(n)$ are all zero if $n'\geq n+n_0$. The lemma is therefore proved for $m=1$.

One now assumes that $m>1$ and that the lemma is proved for the integers $m'<m$; let then $\bbg=(f_1, \ldots, f_{m-1})$ and $\bbh=f_m$.

For every $n>0$, one has (\EGA III~1.1.4.1) an exact sequence:
\begin{equation*}
0 \to \H_0(\bbh^n, \H_i(\bbg^n, N)) \to \H_i(\bbf^n, N) \to \H_1(\bbh^n, \H_{i-1}(\bbg^n, N)) \to 0,
\end{equation*}
and for variable $n$ a projective system of exact sequences. It follows from the induction hypotheses that for $i>0$ the $\H_i(\bbg^n, N)$ form an essentially zero projective system, hence also the $\H_0(\bbh^n, \H_i(\bbg^n, N))$ which one identifies with quotients of $\H_i(\bbg^n, N)$. For the terms on the right one will factor the transition morphisms from $n'$ to $n$ as:
\begin{equation*}
\H_1(\bbh^{n'}, \H_{i-1}(\bbg^{n'}, N)) \to \H_1(\bbh^{n'}, \H_{i-1}(\bbg^{n}, N)) \to \H_1(\bbh^{n}, \H_{i-1}(\bbg^{n}, N)).
\end{equation*}
Since $\H_{i-1}(\bbg^{n}, N)$ is a noetherian module it follows from the case $m=1$ that there exists, for given $n$, $n'>n$ such that the second arrow is zero. One therefore sees that in this projective system of exact sequences, the extreme projective systems are essentially zero, hence the same is true of the middle projective system.

One has therefore proved lemma \Ref{II.11} hence lemma \Ref{II.8} and consequently theorem \Ref{II.6}.

\begin{remarkstar}
One can also obtain theorem \Ref{II.6} by proving the condition of corollary \Ref{II.10} with the aid of the structure theorems for injective modules over a noetherian ring (Matlis, Gabriel).
\end{remarkstar}
